\begin{landscape}
\begingroup
\fontsize{10}{13}\selectfont
\fancyhead[L]{\small NACHLASS UND BRIEFWECHSEL}
\renewcommand{\headrulewidth}{0pt}
\par\medskip
\pdfbookmark[2]{Nachlass VII [2]}{pallas-VII-fragment2}
\begin{center}[2.]\end{center}
\makeatletter
\providecommand{\VIItwoHistoricalNote}[1]{\begingroup
\renewcommand{\thefootnote}{*)}
\long\def\@makefntext##1{\noindent [*)\ ##1}
\footnote{#1}\endgroup}
\makeatother
\par\noindent\hspace*{\parindent}[Um die Berechnung der Störungsbeträge für die einzelnen Oppositionen zu erleichtern, lassen sich in
folgender Weise Tafeln herstellen:
\par\noindent\hspace*{\parindent}Man fasst zunächst die Störungsgleichungen mit den Argumenten
\[
\begin{gathered}
2\text{\mdseries\AstroSymbols ♃}-\text{\mdseries\AstroSymbols ⚴},\quad
\text{\mdseries\AstroSymbols ♃},\quad\text{\mdseries\AstroSymbols ⚴},\quad
\text{\mdseries\AstroSymbols ♃}-2\text{\mdseries\AstroSymbols ⚴},\quad
2\text{\mdseries\AstroSymbols ♃}-3\text{\mdseries\AstroSymbols ⚴}\ \text{etc.}\\
3\text{\mdseries\AstroSymbols ♃}-2\text{\mdseries\AstroSymbols ⚴},\quad
4\text{\mdseries\AstroSymbols ♃}-3\text{\mdseries\AstroSymbols ⚴}\ \text{etc.,}
\end{gathered}
\]
\noindent in eine Gruppe zusammen. Sie sind, wenn man
\[
2\text{\mdseries\AstroSymbols ♃}-\text{\mdseries\AstroSymbols ⚴}=A_1,\qquad
\text{\mdseries\AstroSymbols ⚴}-\text{\mdseries\AstroSymbols ♃}=x
\]
\noindent setzt, von der Form
\[
A_1+nx+G^{(n)},
\]
\noindent wo \(n\) die positiven und negativen ganzen Zahlen, Null eingeschlossen, und \(G^{(n)}\) die zugehörige Constante
bedeutet. Man kann setzen:
\begingroup\small
\[
\begin{aligned}
&g\sin(A_1+G)+g'\sin(A_1+x+G')+g''\sin(A_1+2x+G'')+\cdots\\
&\qquad+h'\sin(A_1-x+H')+h''\sin(A_1-2x+H'')+\cdots\\
{}={}&\sin A_1\,\mathord{.}\,(a+a'\cos x+a''\cos2x+\cdots+b'\sin x+b''\sin2x+\cdots)\\
&+\cos A_1\,\mathord{.}\,(c+c'\cos x+c''\cos2x+\cdots+d'\sin x+d''\sin2x+\cdots),
\end{aligned}
\]
\endgroup
\noindent wo
\begingroup\scriptsize
\[
\begin{array}{l l l}
a=g\cos G,&a'=g'\cos G'+h'\cos H',&a''=g''\cos G''+h''\cos H'',\quad\text{etc.}\\
&b'=-g'\sin G'+h'\sin H',&b''=-g''\sin G''+h''\sin H'',\quad\text{etc.}\\
c=g\sin G,&c'=g'\sin G'+h'\sin H',&c''=g''\sin G''+h''\sin H'',\quad\text{etc.}\\
&d'=g'\cos G'-h'\cos H',&d''=g''\cos G''-h''\cos H'',\quad\text{etc.}
\end{array}
\]
\endgroup
\par\noindent\hspace*{\parindent}Bestimmt man nun \(r\) und \(\varphi\) aus den Gleichungen
\[
\begin{aligned}
\left.\begin{aligned}
a+a'\cos x+a''\cos2x+\cdots\\
{}+b'\sin x+b''\sin2x+\cdots
\end{aligned}\right\}&=P=r\cos\varphi\\[1ex]
\left.\begin{aligned}
c+c'\cos x+c''\cos2x+\cdots\\
{}+d'\sin x+d''\sin2x+\cdots
\end{aligned}\right\}&=Q=r\sin\varphi,
\end{aligned}
\]
\SourceLeaf{574}
\noindent so erhält man für die ganze Gruppe den Ausdruck
\begingroup\small
\[
\left.\begin{aligned}
g\sin(A_1+G)+g'\sin(A_1+x+G')+g''\sin(A_1+2x+G'')+\cdots\\
{}+h'\sin(A_1-x+H')+h''\sin(A_1-2x+H'')+\cdots
\end{aligned}\right\}=r\sin(A_1+\varphi).
\]
\endgroup
\par\noindent\hspace*{\parindent}Die \textls[100]{erste Operation}\VIItwoHistoricalNote{Vgl. den folgenden Brief.]} besteht nun in der Berechnung der \(a^{(n)}\), \(b^{(n)}\), \(c^{(n)}\), \(d^{(n)}\) aus den \(g^{(n)}\), \(h^{(n)}\),
\(G^{(n)}\), \(H^{(n)}\); die \textls[100]{zweite} aus der Berechnung von \(P\) und \(Q\) für die Werthe \(x=0^0\), \(15^0\), \(30^0\), etc., und die
\textls[100]{dritte} aus der von \(r\) und \(\varphi\).
\par\noindent\hspace*{\parindent}In derselben Weise sind die übrigen Störungsgleichungen in Gruppen zusammenzufassen, indem man
beständig
\[
x=\text{\mdseries\AstroSymbols ⚴}-\text{\mdseries\AstroSymbols ♃},
\]
\noindent dagegen
\begingroup\footnotesize
\[
\begin{array}{l|l|l|l}
A_2=3\text{\mdseries\AstroSymbols ♃}-\text{\mdseries\AstroSymbols ⚴}&
A_5=8\text{\mdseries\AstroSymbols ♃}-3\text{\mdseries\AstroSymbols ⚴}&
A_8=13\text{\mdseries\AstroSymbols ♃}-5\text{\mdseries\AstroSymbols ⚴}&
A_{11}=17\text{\mdseries\AstroSymbols ♃}-6\text{\mdseries\AstroSymbols ⚴}\\
A_3=5\text{\mdseries\AstroSymbols ♃}-2\text{\mdseries\AstroSymbols ⚴}&
A_6=10\text{\mdseries\AstroSymbols ♃}-4\text{\mdseries\AstroSymbols ⚴}&
A_9=15\text{\mdseries\AstroSymbols ♃}-6\text{\mdseries\AstroSymbols ⚴}&
A_{12}=20\text{\mdseries\AstroSymbols ♃}-8\text{\mdseries\AstroSymbols ⚴}\\
A_4=6\text{\mdseries\AstroSymbols ♃}-2\text{\mdseries\AstroSymbols ⚴}&
A_7=11\text{\mdseries\AstroSymbols ♃}-4\text{\mdseries\AstroSymbols ⚴}&
A_{10}=16\text{\mdseries\AstroSymbols ♃}-6\text{\mdseries\AstroSymbols ⚴}&
A_{13}=21\text{\mdseries\AstroSymbols ♃}-8\text{\mdseries\AstroSymbols ⚴}
\end{array}
\]
\endgroup
\noindent setzt.
\par\noindent\hspace*{\parindent}Hienach wurde die als Notiz [4.] abgedruckte Tafel entworfen, welche für jedes der \(A\) die Grössen
\(r\) und \(\varphi\) für das Argument \(x\) in Intervallen von \(15^0\) gibt.
\par\noindent\hspace*{\parindent}Die Gruppe der Störungsgleichungen mit den Argumenten \(\text{\mdseries\AstroSymbols ♃}-\text{\mdseries\AstroSymbols ⚴}\), \(2\text{\mdseries\AstroSymbols ♃}-2\text{\mdseries\AstroSymbols ⚴}\), \(3\text{\mdseries\AstroSymbols ♃}-3\text{\mdseries\AstroSymbols ⚴}\) etc. kann direkt
für das Argument \(x\) berechnet werden, ihre Werthe sind im Beginn der Tafel unter der Uberschrift \(A_0=0\)
gegeben.]
\clearpage
\endgroup
\end{landscape}
